\section*{Chapitre \ref{ch:qm:testing-and-multiple-testing} --- Tests et tests multiples}

\begin{solution}{exo:qm:testing-and-multiple-testing:1}
$p = 1 - \Phi(2.5) = 0.0062$~; avec la \omterm{def:qm:testing-and-multiple-testing:fwer}{correction de Bonferroni} sur vingt variantes, $20 \times 0.0062 = 0.124$~: non significatif à 5~\%.
\end{solution}

\begin{solution}{exo:qm:testing-and-multiple-testing:2}
Niveau $1 - \Phi(1.645) = 5\%$~; puissance $\P(\mathcal N(2, 1) > 1.645) = \Phi(0.355) = 0.639$.
\end{solution}

\begin{solution}{exo:qm:testing-and-multiple-testing:3}
Le seuil de Bonferroni vaut 0.0125~: il rejette les deux premières. Holm compare les $p$-valeurs ordonnées à 0.0125, 0.0167, 0.025 et
0.05 tour à tour~: 0.001, 0.012 et 0.02 passent, et 0.3 ne passe pas, si bien que la procédure effectue trois rejets. Les $p$-valeurs ajustées sont, pour
Bonferroni, 0.004, 0.048, 0.08 et 1~; pour Holm, 0.004, 0.036, 0.04 et 0.3.
\end{solution}

\begin{solution}{exo:qm:testing-and-multiple-testing:4}
$T = ((1.645 + 1.282)/0.7)^2 = 17.5$ ans.
\end{solution}

\begin{solution}{exo:qm:testing-and-multiple-testing:5}
Avec une moyenne nulle connue, $\mathrm{LR} = n\ln\hat s^2 - n_1\ln\hat s_1^2 - n_2\ln\hat s_2^2$ avec la variance commune $\hat s^2 = (1 + 1.44)/2 = 1.22$ (en
$\%^2$)~: $1\,000\ln1.22 - 500\ln1.44 = 16.53$. Face à $\chi^2_1$, $p = 4.8 \times 10^{-5}$~: la variance a changé.
\end{solution}

\begin{solution}{exo:qm:testing-and-multiple-testing:6}
$1 - \Phi(2)^{10} = 1 - 0.97725^{10} = 0.206$. La formule donne $(1 - \gamma_E)\Phi^{-1}(0.9) + \gamma_E\Phi^{-1}(1 - 1/(10e)) = 1.57$
(le maximum espéré exact vaut 1.54).
\end{solution}

\begin{solution}{exo:qm:testing-and-multiple-testing:7}
Avec 500 criblages (\texttt{fdr\_correlated()})~: le \omterm{def:qm:testing-and-multiple-testing:fdr}{taux de fausses découvertes} vaut 3.0~\%, sous la borne de 5~\% (dépendance positive), avec
environ 63 découvertes par criblage. Mais la proportion réalisée est volatile~: elle dépasse 10~\% dans 7.8~\% des criblages, lorsque le facteur commun
fait monter ensemble toutes les statistiques nulles. La procédure contrôle une moyenne, et non chaque criblage.
\end{solution}

\begin{solution}{exo:qm:testing-and-multiple-testing:8}
Douze tests à 5~\% produisent en moyenne 0.6 faux positifs~; trois ou plus surviennent avec une probabilité de 2.0~\% si le signal est inutile
partout et les marchés indépendants, si bien que le décompte est une faible indication de quelque chose. Le défaut tient à la sélection~: les trois
gagnants ont été choisis parce que leurs backtests étaient bons, si bien que leur performance en échantillon est biaisée à la hausse (malédiction du vainqueur),
et les marchés ne sont pas indépendants. Présenter le résultat pour la famille, et estimer la performance sur des données non utilisées pour choisir les marchés.
\end{solution}

\begin{solution}{pb:qm:testing-and-multiple-testing:1}
\textbf{1.} Fenêtre de 42 jours~: $t = 3.11$, \omterm{def:qm:estimation:sharpe}{ratio de Sharpe} annualisé 0.98.
\textbf{2.} $p = 1 - \Phi(3.11) = 0.00094$.
\textbf{3.} 33 sur 200.
\textbf{4.} 1.04~: chaque règle porte sur le même historique, si bien que ses épisodes de tendance font monter ensemble toute la famille.
\textbf{5.} 0.95 avec la règle à 43 jours~; 0.28 avec celle à 5 jours et 0.20 avec celle à 204 jours.
\textbf{6.} $200 \times 0.00094 = 0.187$~; valeur critique $\Phi^{-1}(1 - 0.05/200) = 3.48$.
\textbf{7.} 0.187~: Holm multiplie la plus petite $p$-valeur par $m$, comme le fait Bonferroni~; il ne gagne que plus bas dans la liste.
\textbf{8.} 0.076~: la procédure ascendante prend le minimum de $p_{(k)}m/k$ sur les rangs situés au-dessus, et les 33 petites $p$-valeurs rendent ces valeurs petites.
\textbf{9.} $1 - (1 - 0.00094)^{200} = 0.171$.
\textbf{10.} 0.030 à partir de 200\,000 tirages gaussiens avec la corrélation estimée~; la valeur critique à 5~\% du maximum vaut 2.92.
\textbf{11.} $\ln(1 - 0.030)/\ln(1 - 0.00094) = 32$ essais.
\textbf{12.} Trois.
\textbf{13.} Dans 2.6~\% d'entre elles, proche des 3.0~\% simulés (l'erreur d'échantillonnage de 2\,000 historiques vaut 0.4 point).
\textbf{14.} La probabilité que le vrai \omterm{def:qm:estimation:sharpe}{ratio de Sharpe} soit positif vaut 0.999. Déflaté pour 200 essais~: référence à 0.87 par an, \omterm{def:qm:testing-and-multiple-testing:dsr}{ratio de
Sharpe déflaté} de 0.63~; pour 32 essais~: 0.67 et 0.84.
\textbf{15.} À peine~: avec des rendements de Student-$t_4$, la meilleure règle atteint 3.11 dans 2.35~\% de 2\,000 nouveaux historiques~; les statistiques $t$ sont
encore proches de la loi normale après 2\,520 jours.
\textbf{16.} Un test à 5~\% rejette une \omterm{def:qm:testing-and-multiple-testing:test}{hypothèse nulle} vraie dans 5~\% des cas, et la $p$-valeur ajustée de la meilleure règle vaut 3~\%~; de plus, l'historique est le
336\up{e} tiré, une recherche dont la correction n'a pas tenu compte.
\textbf{17.} Un échantillon de test réservé, jamais utilisé dans la recherche, une raison économique à la tendance à cet horizon, une performance nette de coûts, et une
$p$-valeur ajustée qui compte chaque famille testée par l'équipe.
\textbf{18.} Le nombre de variantes et de familles testées, leur corrélation, les $p$-valeurs ajustées et la méthode, la période d'estimation et
ce qui a été fixé avant d'examiner les données.
\textbf{19.} Résultat nommé~: \emph{le meilleur parmi deux cents}~: la meilleure de deux cents règles de momentum corrélées, dont la \omterm{def:qm:testing-and-multiple-testing:test}{p-valeur} naïve vaut $p = 0.00094$, possède une
$p$-valeur ajustée par la statistique du maximum de 0.030 (Bonferroni~: 0.187)~; la recherche équivaut à 32 essais indépendants.
\textbf{20.} Elle protège contre les recherches dont elle a connaissance, et contre rien de ce que le chercheur n'a pas compté.
\end{solution}

\begin{solution}{iq:qm:testing-and-multiple-testing:1}
La probabilité, si l'\omterm{def:qm:testing-and-multiple-testing:test}{hypothèse nulle} est vraie, d'une statistique au moins aussi extrême que celle observée. Ce n'est pas la probabilité que
l'\omterm{def:qm:testing-and-multiple-testing:test}{hypothèse nulle} soit vraie, ni la probabilité que le résultat se répète, ni une mesure de l'ampleur de l'effet.
\iqlookfor{Le sens du conditionnement ($\P(\text{données} \mid H_0)$, et non $\P(H_0 \mid \text{données})$), et le fait que la taille de l'effet demande son propre nombre.}
\end{solution}

\begin{solution}{iq:qm:testing-and-multiple-testing:2}
Probablement pas. Si tout n'était que bruit et indépendant, le meilleur parmi 100 dépasse 2.8 avec une probabilité de $1 - \Phi(2.8)^{100} = 0.23$~;
la corrélation la réduit, mais il faut alors se demander combien de variantes et de choix ont précédé les cent. Corriger pour la recherche (statistique du maximum ou
Bonferroni), puis tester sur un échantillon de test réservé.
\iqlookfor{Le calcul de la chance sur quatre, l'effet de la corrélation et un échantillon de test réservé comme véritable test.}
\end{solution}

\begin{solution}{iq:qm:testing-and-multiple-testing:3}
Bonferroni ou Holm (\omterm{def:qm:testing-and-multiple-testing:fwer}{taux d'erreur familial}) lorsqu'un faux positif est coûteux, comme pour la seule stratégie mise en production~; Benjamini--Hochberg
(\omterm{def:qm:testing-and-multiple-testing:fdr}{taux de fausses découvertes}) lorsqu'on crible de nombreux candidats pour un modèle combiné, où quelques fausses découvertes parmi beaucoup de vraies sont acceptables.
En cas de forte dépendance, utiliser la version par statistique du maximum ou Benjamini--Yekutieli.
\iqlookfor{Adapter le taux d'erreur au coût d'une fausse découverte~; Holm domine Bonferroni.}
\end{solution}

\begin{solution}{iq:qm:testing-and-multiple-testing:4}
L'estimation annuelle a une \omterm{def:qm:estimation:estimator}{erreur type} d'environ $1/\sqrt T$, donc $t = \sqrt T$~: quatre ans pour $t = 2$, et, au niveau de 5~\% avec une puissance de 80~\%,
$(1.645 + 0.842)^2 = 6.2$ ans. Des rendements autocorrélés et des queues épaisses l'allongent~; une recherche l'allonge encore davantage.
\iqlookfor{La règle en $1/\sqrt T$ et le calcul de puissance, et non seulement $t = 2$.}
\end{solution}

\begin{solution}{iq:qm:testing-and-multiple-testing:5}
Utiliser la loi du maximum plutôt que la borne de Boole~: estimer la corrélation des P\&L des variantes et simuler le maximum de
normales corrélées, ou faire un bootstrap conjoint du maximum (\omterm{def:qm:testing-and-multiple-testing:rc}{Reality Check}, SPA, Romano--Wolf step-down). Présenter le nombre effectif
d'essais. La \omterm{def:qm:testing-and-multiple-testing:fwer}{correction de Bonferroni} est valide mais très conservatrice ici.
\iqlookfor{Statistique du maximum plutôt que borne de Boole, rééchantillonnage conjoint qui conserve la dépendance, et comptage des familles autant que des variantes.}
\end{solution}

\begin{solution}{iq:qm:testing-and-multiple-testing:6}
La multiplicité implicite des choix qui dépendent des données, même lorsqu'une seule analyse est menée. En apprentissage automatique~: les choix de variables faits après
avoir examiné les scores de validation, l'arrêt précoce et les hyperparamètres ajustés sur l'ensemble de test, relances avec de nouvelles graines jusqu'à ce que la courbe paraisse bonne,
et un ensemble de test réutilisé au fil de nombreuses itérations de modèle jusqu'à devenir un ensemble d'entraînement.
\iqlookfor{Points de fuite concrets, et la discipline d'un échantillon de test consulté une seule fois.}
\end{solution}

\begin{solution}{iq:qm:testing-and-multiple-testing:7}
Peu. Le test est faible dans les queues, là où les distributions de rendements diffèrent~; avec des paramètres estimés sur les mêmes données, ses
valeurs critiques tabulées sont trop indulgentes (Lilliefors), si bien qu'il rejette encore moins souvent~; et ne pas rejeter n'équivaut pas à accepter. Comparer
directement les quantiles de queue ou les dénombrements de dépassements, avec des valeurs critiques simulées.
\iqlookfor{Faible puissance dans les queues, correction des paramètres estimés, et distinction entre absence de preuve et preuve d'ajustement.}
\end{solution}
